% ======================================================================
\section{Reaction rates handed to the chemistry}
\label{sec:photorates}
% ======================================================================

The chemistry's rate coefficients are computed per cell in \rtn{pdrStep}
(\src{pdr\_mod.f90:\allowbreak{}3180--4580}) by 3D-PDR's
\rtn{CALCULATE\_REACTION\_RATES} (\src{pdr/\allowbreak{}calc\_reac\_rates.F90}), called at
\src{pdr\_mod.f90:\allowbreak{}3692} with MOCASSIN's inputs, and then partly overwritten
with rates that MOCASSIN or the level model own. Table~\ref{tab:rateorigin}
summarises where each radiation-driven rate comes from; reaction numbers are
those of \src{dataPDR/\allowbreak{}rates\_reduced.d}.

\begin{longtable}{@{}L{0.28\linewidth}L{0.15\linewidth}L{0.37\linewidth}L{0.09\linewidth}@{}}
\caption{Origin of the radiation- and cosmic-ray-driven rates in the 3D build.}\label{tab:rateorigin}\\
\toprule
Reaction(s) & Number(s) & Rate used & Section \\
\midrule
\endfirsthead
\toprule
Reaction(s) & Number(s) & Rate used & Section \\
\midrule
\endhead
\Htwo\,+\,$h\nu\to$\,2H & 329 & the level model's $k_{\mathrm{Sol}}$ & \ref{sec:h2rate} \\
CO\,+\,$h\nu\to$\,C\,+\,O & 290 & $\alpha_{\mathrm{CO}}\sum_r\chi_rf_r\Theta$ \citep{vanDishoeckBlack1988} & \ref{sec:co} \\
24 other photoreactions & 266--268, 270--288, 291--292 & $\alpha\chi$ (Draine spectral shape) & \ref{sec:generic} \\
C, Mg\,+\,$h\nu\to$\,ion\,+\,$e^-$ & 269, 289 & set to zero & \ref{sec:linked} \\
ionisation of H, He, C, O, Mg & 293, 297, 298, 299, 312 & cosmic-ray (or, Mg, induced-photon) rate plus MOCASSIN's photoionisation rate & \ref{sec:linked} \\
\Htwo\,+\,CRP\,$\to$\,\Htwo$^+$\,+\,$e^-$ & 296 & cosmic-ray rate plus MOCASSIN's \Htwo{} photoionisation & \ref{sec:linked} \\
X$^+$\,+\,$e^-\to$\,X\,+\,$h\nu$ (X = H, He, C, O, Mg) & 242--246, 248, 249 & MOCASSIN's total recombination coefficient & \ref{sec:linked} \\
cosmic-ray ionisation & 293--300 & $\alpha\,\zeta/\zeta_0$ & \ref{sec:cr} \\
cosmic-ray-induced photoreactions & 301--315 & $\alpha(\zeta/\zeta_0)(T/300)^\beta\gamma/(1-\omega)$ & \ref{sec:cr} \\
charge transfer with H, \Hp & 156--157, 161--163, 167 & \rtn{pdrCtCoef} & \ref{sec:ct} \\
H\,+\,H\,(grain)\,$\to$\,\Htwo & 328 & \citet{CazauxTielens2002} & \ref{sec:formation} \\
\bottomrule
\end{longtable}

\subsection{The field as a set of rays}
\label{sec:rays}

3D-PDR evaluates each photoreaction as a sum over rays, each ray $r$ carrying
a field strength $\chi_r$ (in Draine units), a visual extinction $A_{V,r}$ and
shielding columns. In the 3D build the bridge builds
$n_{\mathrm{chem}} = 1 + N_{\mathrm{site}}$ (+1 fallback) rays
(\src{pdr\_mod.f90:\allowbreak{}3383--3455}, \dlog{D-348}, \dlog{D-350}), which
represent the local MC field split by the paths its photons took
(Sections~\ref{sec:split} and \ref{sec:sites}):

\begin{longtable}{@{}L{0.17\linewidth}L{0.20\linewidth}L{0.10\linewidth}L{0.17\linewidth}L{0.24\linewidth}@{}}
\caption{The rays handed to the chemistry in the 3D build.}\label{tab:rays}\\
\toprule
Ray & $\chi_r$ & $A_{V,r}$ & LW factor & Columns $N(\mathrm{H_2}),N(\mathrm{CO})$ \\
\midrule
\endfirsthead
\toprule
Ray & $\chi_r$ & $A_{V,r}$ & LW factor & Columns $N(\mathrm{H_2}),N(\mathrm{CO})$ \\
\midrule
\endhead
0 (direct) & $\chiD$ & 0 & $f_{\mathrm{dir}}$ & source-ward, Eq.~(\ref{eq:avradial}) \\
site $j$ & $\chiF\,w_j/\sum w$ & 0 & $f_{\mathrm{dif}}$ & two-leg, Eq.~(\ref{eq:twoleg}) \\
fallback & $\chiF$ & 0 & $f_{\mathrm{dif}}$ & source-ward (only when $\chiF>0$ and no site exists) \\
\bottomrule
\end{longtable}

$A_{V,r}=0$ on every ray (\src{:\allowbreak{}3439}) because the MC field has already been
attenuated by the dust; applying 3D-PDR's $e^{-\gamma A_V}$ again would count
the dust twice (\dlog{D-343}, \dlog{D-350}). The C\,\textsc{i} column handed
over is zero (\src{:\allowbreak{}3627}) because MOCASSIN's opacity already contains the
C\,\textsc{i} photoionisation continuum; the HD and S columns are zero
(neither species is in the network). If a source-ward column is missing
(negative), it falls back to $(\AV/(\AV/N_{\mathrm H}))\,x_X$
(\src{:\allowbreak{}3458}, \src{:\allowbreak{}3603--3605}); zero is a valid column and is kept
(\dlog{D-088}). The columns exclude the target cell (Section~\ref{sec:avradial});
only the level model averages the shielding through the cell's own gas
(Section~\ref{sec:ownavg}).

\subsection{Generic photoreactions}
\label{sec:generic}

For every reaction $X+\mathrm{PHOTON}\to\ldots$ without a dedicated treatment,
3D-PDR's rate is $k=\sum_r\alpha\,\chi_r\,e^{-\gamma A_{V,r}}$ [s$^{-1}$]
(\src{pdr/\allowbreak{}calc\_reac\_rates.F90:\allowbreak{}321--347}), with $\alpha$ the rate in the
unattenuated Draine field and $\gamma$ the dust-attenuation exponent of the UMIST
format; with $A_{V,r}=0$,
\begin{equation}
  k = \alpha\,(\chiD+\chiF) = \alpha\,\chi ,
  \label{eq:generic}
\end{equation}
the coefficient $\alpha$ of the rate file times the local MC field in the
6--13.6\,eV band. This applies to the photodissociation and photoionisation of
H$_2^+$, H$_3^+$, CH, CH$^+$, CH$_2$, CH$_2^+$, CH$_3$, CH$_3^+$, CH$_4$, OH,
OH$^+$, H$_2$O and O$_2$ (24 of the network's 28 photoreactions). The
coefficients $\alpha$ were computed for the Draine spectral shape, while the
MC field reddens with depth: its LW share falls from 0.30 to 0.008 through
\texttt{hiipdr}'s PDR (\dlog{D-349}). Reactions whose cross sections lie mainly
above 11.2\,eV are therefore overestimated at depth, and those driven by the
6--11\,eV part underestimated; none of these reactions is shielded by \Htwo{}
or CO lines. Rates computed from cross sections folded over the actual field
would remove the assumption; \citet{Heays2017} provide cross sections for 102
species for that purpose, and find that their rates for the interstellar field
agree with earlier compilations to about 30\% in general.

\subsection{\texorpdfstring{\Htwo}{H2} photodissociation: the level model's rate}
\label{sec:h2rate}

3D-PDR first computes its own rate on the rays (\rtn{H2PDRATE}, \rtn{H2SHIELD1};
\src{pdr/\allowbreak{}shield.F90:\allowbreak{}27--55, 231--281}; \src{pdr/\allowbreak{}calc\_reac\_rates.F90:\allowbreak{}259--268}),
\begin{equation}
  k^{\mathrm{net}}_{\mathrm{H_2}} = \sum_r \alpha_{\mathrm{H_2}}\,\chi_r f_r\,
     J_{\mathrm{FGK}}(N_{\mathrm{H_2},r})\,\rtn{SCATTER}(0,1000\,\text{\AA}),
  \qquad \alpha_{\mathrm{H_2}}=5.18\times10^{-11}\unit{s^{-1}},
\end{equation}
where $J_{\mathrm{FGK}}=J_D+J_R$ is the single-line self-shielding function of
\citet[their eqs.~A7--A9]{Federman1979} for one representative line at
1000\,\AA{} with oscillator strength 0.01, half the \Htwo{} in its lower level, a
Doppler width $v_{\mathrm{turb}}/\lambda$ and a radiative width of $8\times10^7$\,s$^{-1}$
(\src{pdr/\allowbreak{}shield.F90:\allowbreak{}46--51, 249--252}). The bridge then replaces it with
the level model's Solomon rate (\dlog{H-012}):
\begin{equation}
  k(\mathrm{H_2}+h\nu\to2\mathrm H) = k_{\mathrm{Sol}}
  \qquad(\src{pdr\_mod.f90:\allowbreak{}3736}).
\end{equation}
The same rate enters 3D-PDR's heating routine, so the chemistry and the
heating use one \Htwo{} photodissociation rate. The network's own rate is kept
as a diagnostic (\rtn{h2SolChem}, on the same field split and columns). On
\texttt{hiipdr} the level model's rate is 2.4--2.6 times the network's in the
ionised gas, 1.5--1.8 times in cold unshielded gas, and 0.6--0.7 and 0.4--0.6
times behind source-ward columns of $10^{15}$--$10^{17}$\,cm$^{-2}$ and more
(\dlog{H-007}). Three things separate them: the level model's rate in the
unshielded Draine field is 1.23 times $\alpha_{\mathrm{H_2}}$ and the stellar field is
harder; in warm gas the molecules sit in excited levels that dissociate
faster; and the network shields one line of fixed radiative width at the near
face of the cell, whereas the level model averages every line through the cell
and includes overlap.

\subsection{CO photodissociation}
\label{sec:co}

\rtn{COPDRATE} (\src{pdr/\allowbreak{}shield.F90:\allowbreak{}83--106}; called at
\src{pdr/\allowbreak{}calc\_reac\_rates.F90:\allowbreak{}283--292}):
\begin{equation}
  k_{\mathrm{CO}} = \alpha_{\mathrm{CO}}\sum_r \chi_r\,f_r\;
     \Theta\bigl(N_{\mathrm{CO},r},N_{\mathrm{H_2},r}\bigr)\;\rtn{SCATTER}(0,\bar\lambda),
  \qquad \alpha_{\mathrm{CO}} = 2.0\times10^{-10}\unit{s^{-1}},
\end{equation}
with $\alpha_{\mathrm{CO}}$ from the rate file (reaction 290). $\Theta$ is the combined
CO self-shielding and \Htwo{} screening factor for $^{12}$CO of
\citet[their Table~5]{vanDishoeckBlack1988}, interpolated by a two-dimensional
cubic spline in $\log_{10}\Theta$ over $(\log_{10}N_{\mathrm{CO}},\log_{10}N_{\mathrm{H_2}})$ and
clamped at the table edges (\rtn{COSHIELD}, \src{pdr/\allowbreak{}shield.F90:\allowbreak{}375--416}).
The table spans $N(\mathrm{CO})=10^{13}$--$10^{19}$ and
$N(\mathrm{H_2})=10^{19}$--$10^{23}$\,cm$^{-2}$ plus the unshielded entries, which 3D-PDR
places at $10^{12}$ and $10^{18}$\,cm$^{-2}$ (\src{pdr/\allowbreak{}modules.F90:\allowbreak{}353--356},
\src{pdr/\allowbreak{}allocations.F90:\allowbreak{}9--14}); smaller columns are treated as these.
$\bar\lambda$ is the mean wavelength of the 33 dissociating bands weighted by
their contributions to the total rate at the given columns, at which the dust
attenuation is evaluated (\rtn{LBAR}, their eq.~4,
\src{pdr/\allowbreak{}shield.F90:\allowbreak{}570--599}):
\begin{equation}
  \bar\lambda = \bigl[(5675-200.6\,w) - (571.6-24.09\,w)\,u + (18.22-0.7664\,w)\,u^2\bigr]\ \text{\AA},
  \label{eq:lbar}
\end{equation}
with $u=\log_{10}N_{\mathrm{CO}}$ and $w=\log_{10}N_{\mathrm{H_2}}$ [cm$^{-2}$],
clamped to 913.6--1076.1\,\AA. With $A_V=0$, $\rtn{SCATTER}=1$ and $\bar\lambda$ has
no effect. The window factor $f_r$ corrects for the spectral shape of the MC
field in the LW window (Section~\ref{sec:window}); CO is the one photoreaction
in this branch that uses it. The columns are those of the chemistry's own
$n(\mathrm{H_2})$ and $n(\mathrm{CO})$, on the source-ward or two-leg paths, and
exclude the target cell, so CO is under-shielded in the cells of its own
transition (Section~\ref{sec:limitations}).

The shielding functions of \citet{vanDishoeckBlack1988} were computed for
the unmodified Draine field, cool gas ($T\lesssim100$\,K, CO excitation
temperature 10\,K), Doppler parameters of 1\,km\,s$^{-1}$ for CO and 3\,km\,s$^{-1}$ for \Htwo{}
and a fixed H\,\textsc{i} column of $5\times10^{20}$\,cm$^{-2}$ (notes to their Table~5).
\citet{Visser2009} recomputed them from newer spectroscopic data and find an
unattenuated rate of $2.6\times10^{-10}$\,s$^{-1}$, 30\% above the 2.0$\times10^{-10}$ used
here, and depth-dependent rates that change by up to a factor of three when
the excitation temperature or the Doppler width is raised, for temperatures
above 100\,K. The warm CO of the illuminated PDR is therefore outside the
conditions of the table, and the \Htwo{} screening uses the chemistry's total
\Htwo{} column rather than the level model's line-by-line transmission
(Section~\ref{sec:limitations}).

\subsection{Photoionisation and recombination of the linked atoms}
\label{sec:linked}

Five elements are \emph{linked}: their neutral and first ion are species of
the network and stages of MOCASSIN (H, He, C, O, Mg;
\src{pdr\_mod.f90:\allowbreak{}553--564}). For them the network's own photoionisation is
replaced by MOCASSIN's, computed on the full MC field with the real cross
sections (\dlog{D-358}; \src{pdr\_mod.f90:\allowbreak{}3706--3722}; slots found by structure
at \src{:\allowbreak{}743--801}):
\begin{itemize}
\item the network reaction $X+\mathrm{PHOTON}\to X^++e^-$ is set to zero where it
  exists (C, reaction 269; Mg, 289; the network has none for H, He and O);
\item MOCASSIN's photoionisation rate from the ground state
  \begin{equation}
    \Gamma_X = \sum_k \frac{\Jste_k\,\sigma_X(\nu_k)}{h\nu_k}\quad[\mathrm{s^{-1}}]
    \label{eq:gammaX}
  \end{equation}
  (\src{update\_mod.f90:\allowbreak{}303--311}; outer-shell cross sections of
  \citealt{Verner1996} and \citealt{VernerYakovlev1995}, \src{ph\_mod.f90:\allowbreak{}473--476})
  is added to the reaction's first-order cosmic-ray slot $X+\mathrm{CRP}\to X^++e^-$
  (H, He, C, O: reactions 293, 297, 298, 299), on top of the cosmic rays; Mg has
  no such reaction and its cosmic-ray-induced photoionisation (312) takes the
  rate instead;
\item for hydrogen, MOCASSIN's collisional ionisation term
  $2.75\times10^{-16}\,T^{3/2}(157\,894/T+2)\,e^{-157\,894/T}$ is added as well
  (\src{update\_mod.f90:\allowbreak{}2620--2630}, stored at \src{:\allowbreak{}2859}). It has the form of a
  rate coefficient [cm$^3$\,s$^{-1}$]: at $10^4$\,K it is $6.8\times10^{-16}$, the
  magnitude of the H\,\textsc{i} collisional ionisation coefficient, but it is added
  to the photoionisation rate [s$^{-1}$] without the factor $n_e$, here and in
  MOCASSIN's own ionisation balance (\src{:\allowbreak{}2769--2770}), so collisional
  ionisation is underestimated by the factor $n_e$. With that factor it would
  reach 10\% of a photoionisation rate of $10^{-8}$\,s$^{-1}$ at about
  $1.5\times10^4$\,K for $n_e=10^4$\,cm$^{-3}$; the omission matters in hot, dense
  ionised gas and not in the PDR. The source comment
  attributes the formula to \citet{DrakeUlrich1980}; that paper takes its
  collisional ionisation rates from \citet{Johnson1972} and tabulates
  $4.88\times10^{-16}$\,cm$^3$\,s$^{-1}$ for the ground level at $10^4$\,K (their
  Table~3), so the formula's origin is not established
  (Section~\ref{sec:limitations});
\item the radiative recombination $X^++e^-\to X+\mathrm{PHOTON}$ takes
  MOCASSIN's total recombination coefficient (radiative, from the fits of
  \citealt{VernerFerland1996} and the others listed at
  \src{update\_mod.f90:\allowbreak{}4101--4115}, plus dielectronic) in the first of its
  temperature ranges, the other ranges set to zero;
\item \Htwo{} photoionisation, $\mathrm{H_2}+h\nu\to\mathrm{H_2^+}+e^-$, is
  added to the slot $\mathrm{H_2}+\mathrm{CRP}\to\mathrm{H_2^+}+e^-$ (296) with
  $\Gamma_{\mathrm{H_2}}=\sum_{h\nu_k\ge15.4\,\mathrm{eV}}\Jste_k\sigma_{\mathrm{H_2}}(\nu_k)/h\nu_k$
  and the fit of \citet{Yan1998}, as revised in their erratum \citep{Yan2001},
  to the cross section (\src{update\_mod.f90:\allowbreak{}2863--2877}; \rtn{yanH2xSec},
  \src{ionization\_mod.f90:\allowbreak{}385--411}).
\end{itemize}
These rates are stored by \rtn{updateCell} at every trial temperature of its
search, so the chemistry reads those of the temperature it settled on
(\src{update\_mod.f90:\allowbreak{}2838--2878}). In particular, \textbf{carbon is
photoionised at MOCASSIN's rate}: the MC field contains dust and C\,\textsc{i}
continuum absorption but no \Htwo{} line absorption, so the C\,\textsc{i} rate is
not reduced by the \Htwo{} lines that overlap the C\,\textsc{i} continuum
(Section~\ref{sec:limitations}). The C\,\textsc{i} treatment of
\citet{KampBertoldi2000} in 3D-PDR (\rtn{CIPDRATE},
\src{pdr/\allowbreak{}shield.F90:\allowbreak{}132--157}),
$k=\alpha\chi\exp[-\gamma\AV-1.1\times10^{-17}N_{\mathrm C}-0.9\,T^{0.27}(N_{\mathrm{H_2}}/1.59\times10^{21})^{0.45}]$,
is computed with $\AV=0$ and $N_{\mathrm C}=0$, i.e.\ with its \Htwo{} term only,
and then discarded. (Equation~8 of \citealt{KampBertoldi2000}, which the code
cites, normalises the \Htwo{} term at $10^{22}$\,cm$^{-2}$, not $1.59\times10^{21}$; the
vendored term is 2.3 times larger at a given column.) The ionisation balance of
the linked elements is moreover imposed from MOCASSIN before and after the
solve (Section~\ref{sec:impose}).

\subsection{Cosmic rays and cosmic-ray-induced photons}
\label{sec:cr}

The cosmic-ray ionisation rate is the input \kw{pdrZeta} (default
$10^{-16}$\,s$^{-1}$, \src{common\_mod.f90:\allowbreak{}902}), stored normalised to the
UMIST reference rate $\zeta_0=1.3\times10^{-17}$\,s$^{-1}$ (\src{pdr\_mod.f90:\allowbreak{}360}),
and is the same in every cell (3D-PDR's attenuation option \kw{CRATTENUATION}
is not compiled). Direct cosmic-ray reactions have $k=\alpha\,\zeta/\zeta_0$, with
$\alpha$ the rate at $\zeta_0$ (\src{pdr/\allowbreak{}calc\_reac\_rates.F90:\allowbreak{}353--370}). Reactions
driven by the ultraviolet photons that the secondary electrons of
cosmic-ray ionisation excite in \Htwo{} \citep{PrasadTarafdar1983} have
\begin{equation}
  k = \alpha\,\frac{\zeta}{\zeta_0}\left(\frac{T}{300\,\mathrm K}\right)^{\beta}\frac{\gamma}{1-\omega}
  \quad[\mathrm{s^{-1}}]
  \label{eq:crphot}
\end{equation}
(\src{pdr/\allowbreak{}calc\_reac\_rates.F90:\allowbreak{}381--401}), the UMIST form of the rate of
\citet[their eq.~7]{Gredel1989}, $\zeta\,p_M/(1-\omega)$ per molecule: $\alpha=\zeta_0$ for
all 15 such reactions of the network, $\gamma=p_M$ is the efficiency of the
photoreaction per ionisation, and $\omega$ the grain albedo, which enters because
the induced photons are removed by grain absorption, $\sigma_g(1-\omega)$, rather
than by extinction. That approximation holds when grains dominate the
absorption of the induced photons, which \citet{Gredel1989} note is the case
for all molecules except CO; the temperature factor is used only by CO
($\beta=1.17$, reaction 313). Here $\omega$ is MOCASSIN's grain albedo in the LW
band (Section~\ref{sec:dust}). The network has 8 \kw{CRP} and 15 \kw{CRPHOT}
reactions. \citet{Heays2017} find newly calculated cosmic-ray-induced rates
that differ from the standard values by up to a factor of five.

\subsection{Charge transfer with hydrogen}
\label{sec:ct}

The charge-transfer reactions of the linked elements with H and \Hp{}
(reactions 156--157, He$^+$\,+\,H; 161--163, O$^+$\,+\,H and \Hp\,+\,O; 167,
\Hp\,+\,Mg; the network has none for carbon) use one coefficient per reaction
in both the chemistry and MOCASSIN's ionisation balance (\rtn{pdrCtCoef},
\src{common\_mod.f90:\allowbreak{}1383--1436}; \src{pdr\_mod.f90:\allowbreak{}3742--3747}; \dlog{D-403}): the
network's $\alpha(T/300)^\beta e^{-\gamma/T}$ in its temperature ranges, replaced by the
fits of \citet{KingdonFerland1996}, $k=a\,t_4^{\,b}(1+c\,e^{d\,t_4})$ with $t_4=T/10^4$\,K,
inside their ranges for He$^+$\,+\,H (6000--$10^5$\,K) and \Hp\,+\,Mg
(5000--30\,000\,K) (\src{pdr\_mod.f90:\allowbreak{}848--855}). For oxygen the network's rates
are fits to \citet{Stancil1999} and are kept. The coefficients are evaluated
at the cell's own $T_e$, not at the clamped temperature of the other rates.

\subsection{\texorpdfstring{\Htwo}{H2} formation on grains}
\label{sec:formation}

The reaction $\mathrm H+\mathrm H\to\mathrm{H_2}$ on grains (reaction 328) takes the
rate of \citet[their eqs.~15--18]{CazauxTielens2002}, whose model
\citet{CazauxTielens2004} develop in full (\rtn{H2\_FORMATION\_RATE},
\src{pdr/\allowbreak{}h2\_form.F90:\allowbreak{}9--134}, flag \kw{H2FORM\_CT02}; called at
\src{pdr/\allowbreak{}calc\_reac\_rates.F90:\allowbreak{}116}):
\begin{gather}
  R_{\mathrm f} = \tfrac12\,v_{\mathrm{th}}\,
     \bigl(\sigma^{\mathrm{sil}}_{\mathrm H}\,\epsilon_{\mathrm{sil}}+\sigma^{\mathrm{gra}}_{\mathrm H}\,\epsilon_{\mathrm{gra}}\bigr)\,
     S_{\mathrm H}(T,T_d)\quad[\mathrm{cm^3\,s^{-1}}],
  \label{eq:rform}\\
  v_{\mathrm{th}} = 1.45\times10^{5}\sqrt{T/100\,\mathrm K}\unit{cm\,s^{-1}},\qquad
  S_{\mathrm H} = \bigl[1+0.04\sqrt{T+T_d}+0.002\,T+8\times10^{-6}\,T^2\bigr]^{-1},\\
  \epsilon = \frac{\xi}{1+F_1+F_2},\quad
  F_1 = \frac{\mu F}{2\nu_{\mathrm{H_2}}e^{-E_{\mathrm{H_2}}/T_d}},\quad
  F_2 = \frac{\bigl(1+\sqrt{(E_{\mathrm{HC}}-E_{\mathrm S})/(E_{\mathrm{HP}}-E_{\mathrm S})}\bigr)^2}{4}\,e^{-E_{\mathrm S}/T_d},\\
  \xi = \Bigl[1+\frac{\nu_{\mathrm{HC}}}{2F}\,e^{-1.5E_{\mathrm{HC}}/T_d}
        \bigl(1+\sqrt{(E_{\mathrm{HC}}-E_{\mathrm S})/(E_{\mathrm{HP}}-E_{\mathrm S})}\bigr)^2\Bigr]^{-1},
\end{gather}
with temperatures in K. $v_{\mathrm{th}}$ is the mean speed of H atoms, the factor
$\frac12$ makes one molecule of two atoms, and $\sigma_{\mathrm H}$ [cm$^2$] is the
geometric grain cross section per H nucleus. $S_{\mathrm H}$ is the sticking
coefficient of \citet[eq.~3.7]{HollenbachMcKee1979}, which they take from
calculations by Burke and Hollenbach; it is 0.58, 0.33 and 0.08 at $T=100$,
300 and 1000\,K (for $T_d=20$\,K). The recombination efficiency $\epsilon$ is
limited by three processes \citep{CazauxTielens2002}: $F_1$, newly formed
molecules that stay on a very cold surface; $F_2$, the evaporation of
physisorbed atoms before they meet a chemisorbed one, which controls the
decline above a few tens of K; and $\xi$, the evaporation of chemisorbed atoms at
several hundred K. The constants are $F=10^{-10}$ monolayers\,s$^{-1}$ (the
accretion rate), $\mu=0.005$, $\nu_{\mathrm{H_2}}=3\times10^{12}$\,s$^{-1}$,
$\nu_{\mathrm{HC}}=1.3\times10^{13}$\,s$^{-1}$, $E_{\mathrm{HC}}=3\times10^4$\,K, and
$(E_{\mathrm S},E_{\mathrm{H_2}},E_{\mathrm{HP}})=(110,320,450)$\,K for silicate and
$(260,520,800)$\,K for graphite. The functional form and $\mu$ are those of
\citet{CazauxTielens2002}; the per-material energies are those tabulated by
\citet[Table~C.1]{Rollig2013} from the work of Cazaux and Tielens, whereas
\citet[Table~1]{CazauxTielens2002} list $(E_{\mathrm S},E_{\mathrm{HP}},E_{\mathrm{HC}})=(200,600,10^4)$\,K
for their silicate surface, the values Cloudy uses; \citet{Rollig2013} also
set the $F_1$ term to zero. With the vendored parameters
(Equation~\ref{eq:rform} evaluated for this document) $\epsilon_{\mathrm{sil}}$ is
close to 1 between 6 and 15\,K, 0.90 at 20\,K, 0.25 at 50\,K and 0.10 at 100\,K, and
$\epsilon_{\mathrm{gra}}$ close to 1 between 10 and 40\,K, 0.43 at 100\,K and 0.12 at 300\,K;
both fall to zero below about 6 and 9\,K ($F_1$) and above about 800\,K ($\xi$).
The silicate energies are the main reason why the silicate coefficient is
0.50--0.91 of Cloudy's on \texttt{hiipdr} (\dlog{L-33}, \dlog{H-022}).

The grain areas per H $\sigma^{\mathrm{sil}}_{\mathrm H}$, $\sigma^{\mathrm{gra}}_{\mathrm H}$ are
MOCASSIN's (Section~\ref{sec:grainarea}), handed over per cell at
\src{pdr\_mod.f90:\allowbreak{}3295--3307}; $T_d$ is MOCASSIN's mean grain temperature
(Section~\ref{sec:Td}), floored at $T_{\mathrm{CMB}}$; the metallicity factor is 1
(\src{pdr\_mod.f90:\allowbreak{}347}). The formation rate per unit volume is
$R_{\mathrm f}\,n_{\mathrm H}\,n(\mathrm H)$ [cm$^{-3}$\,s$^{-1}$]. The same function, evaluated
with one material's area at a time, weights the two formation distributions of
the level model (Section~\ref{sec:formdist}).
